\documentclass[11pt]{article}
\usepackage{amsmath,amssymb,amsthm}
\usepackage{algorithm}
\usepackage{algorithmic}
\usepackage{geometry}
\geometry{margin=1in}
\setlength{\parskip}{0.5em}
\setlength{\parindent}{0pt}
\newtheorem{theorem}{Theorem}
\newtheorem{lemma}{Lemma}
\newtheorem{assumption}{Assumption}

\begin{document}

%%%%%%%%%%%%%%%%%%%%%%%%%%%%%%%%%%%%%%%%%%%%%%%%%%%%%%%%%%%%%%%%%%%%%%%%%%%%
\section*{Algorithm Name}
\textbf{Federated Local Stochastic Gradient Descent (FedAvg)}  
\(K\) clients perform \(E\) local SGD steps with stepsize \(\eta>0\) and
periodically synchronize by averaging their models.

%%%%%%%%%%%%%%%%%%%%%%%%%%%%%%%%%%%%%%%%%%%%%%%%%%%%%%%%%%%%%%%%%%%%%%%%%%%%
\section*{Mathematical Formulation}

\noindent\textbf{Notation}
\begin{itemize}
  \item Global model at communication round \(t\): \(w_t\in\mathbb{R}^d\).
  \item Local model of client \(k\) after the \(e\)-th local step in round \(t\):
        \(w_{t}^{(k,e)}\) with \(w_{t}^{(k,0)}=w_t\).
  \item Stochastic gradient on client \(k\) at step \((t,e)\):
        \(g_{t}^{(k,e)}:=\nabla f_k\!\left(w_{t}^{(k,e)};\xi_{t}^{(k,e)}\right)\).
  \item Local objective on client \(k\): \(F_k(w)=\mathbb{E}_{\xi}[f_k(w;\xi)]\).
  \item Global objective: \(F(w)=\frac1K\sum_{k=1}^{K}F_k(w)\).
  \item Client drift (difference between true local gradient at the
        global point and the average of the gradients seen locally):
        \[
        \delta_t^{(k)}\;:=\;\frac1E\sum_{e=0}^{E-1}
        \bigl(\nabla F_k(w_t)-\nabla F_k(w_{t}^{(k,e)})\bigr).
        \]
\end{itemize}

\noindent\textbf{Local updates (for each client \(k\))}  
\[
\boxed{
\begin{aligned}
w_{t}^{(k,0)} &= w_t,\\
w_{t}^{(k,e+1)} &= w_{t}^{(k,e)} - \eta\, g_{t}^{(k,e)},
\qquad e=0,\dots,E-1 .
\end{aligned}}
\tag{1}
\]

\noindent\textbf{Server aggregation}  
After the \(E\) local steps each client sends \(w_{t}^{(k,E)}\) to the
server, which forms the next global model by simple averaging:
\[
\boxed{ \; w_{t+1}= \frac1K\sum_{k=1}^{K} w_{t}^{(k,E)}\; } .
\tag{2}
\]

The algorithm repeats (1)–(2) for \(t=0,1,\dots,T-1\).

%%%%%%%%%%%%%%%%%%%%%%%%%%%%%%%%%%%%%%%%%%%%%%%%%%%%%%%%%%%%%%%%%%%%%%%%%%%%
\section*{Pseudocode}
\begin{algorithm}[H]
\caption{Federated Local SGD (FedAvg)}
\begin{algorithmic}[1]
\REQUIRE Number of clients \(K\), local steps \(E\), stepsize \(\eta>0\),
          total communication rounds \(T\)
\STATE Initialize global model \(w_0\)
\FOR{communication round \(t=0\) \TO \(T-1\)}
    \FOR{each client \(k=1\) \TO \(K\) \textbf{in parallel}}
        \STATE \(w_{t}^{(k,0)} \gets w_t\)
        \FOR{local step \(e=0\) \TO \(E-1\)}
            \STATE Sample minibatch \(\xi_{t}^{(k,e)}\)
            \STATE Compute stochastic gradient 
                  \(g_{t}^{(k,e)} \gets \nabla f_k\!\bigl(w_{t}^{(k,e)};\xi_{t}^{(k,e)}\bigr)\)
            \STATE Update local model 
                  \(w_{t}^{(k,e+1)} \gets w_{t}^{(k,e)} - \eta\, g_{t}^{(k,e)}\)
        \ENDFOR
        \STATE Send \(w_{t}^{(k,E)}\) to the server
    \ENDFOR
    \STATE Server aggregates: \(w_{t+1}\gets \frac1K\sum_{k=1}^{K} w_{t}^{(k,E)}\)
\ENDFOR
\RETURN Final model \(w_T\) (or average \(\bar w_T:=\frac1T\sum_{t=0}^{T-1} w_t\))
\end{algorithmic}
\end{algorithm}

%%%%%%%%%%%%%%%%%%%%%%%%%%%%%%%%%%%%%%%%%%%%%%%%%%%%%%%%%%%%%%%%%%%%%%%%%%%%
\section*{Dry Run Example}
Consider a scalar quadratic loss
\[
f(w)=(w-3)^2 ,\qquad \nabla f(w)=2(w-3).
\]
Take a single client (\(K=1\)), \(E=1\) (no local drift), stepsize
\(\eta=0.1\), and start from \(w_0=0\).

\begin{center}
\begin{tabular}{c|c|c|c}
Iteration \(t\) & Gradient \(g_t\) & Update \(w_{t+1}=w_t-\eta g_t\) & Objective \(f(w_{t+1})\)\\\hline
0 & \(2(0-3)=-6\) & \(0-0.1(-6)=0.6\) & \((0.6-3)^2=5.76\)\\
1 & \(2(0.6-3)=-4.8\) & \(0.6-0.1(-4.8)=1.08\) & \((1.08-3)^2=3.6864\)\\
2 & \(2(1.08-3)=-3.84\) & \(1.08-0.1(-3.84)=1.464\) & \((1.464-3)^2=2.3689\)
\end{tabular}
\end{center}
The objective decreases monotonically, illustrating the local SGD step.

%%%%%%%%%%%%%%%%%%%%%%%%%%%%%%%%%%%%%%%%%%%%%%%%%%%%%%%%%%%%%%%%%%%%%%%%%%%%
\section*{Assumptions}
\begin{assumption}[Convexity]\label{as:convex}
Each local objective \(F_k\) is convex:
\[
F_k(\theta w+(1-\theta)u)\le \theta F_k(w)+(1-\theta)F_k(u),
\qquad\forall w,u\in\mathbb{R}^d,\ \theta\in[0,1].
\]
\end{assumption}
\begin{assumption}[L‑smoothness]\label{as:smooth}
There exists \(L>0\) such that for any \(w,u\):
\[
\|\nabla F_k(w)-\nabla F_k(u)\|\le L\|w-u\|,
\qquad\forall k=1,\dots,K.
\tag{3}
\]
\end{assumption}
\begin{assumption}[Unbiased stochastic gradients and bounded variance]\label{as:unbiased}
For every client \(k\) and local step \((t,e)\),
\[
\mathbb{E}\!\bigl[g_{t}^{(k,e)}\mid w_{t}^{(k,e)}\bigr]
    =\nabla F_k\!\bigl(w_{t}^{(k,e)}\bigr),\qquad
\mathbb{E}\!\bigl\|g_{t}^{(k,e)}-\nabla F_k(w_{t}^{(k,e)})\bigr\|^{2}
    \le \sigma^{2}.
\tag{4}
\]
\end{assumption}
\begin{assumption}[Bounded client drift]\label{as:drift}
Define the drift \(\delta_t^{(k)}\) as above. Assume
\[
\mathbb{E}\bigl\|\delta_t^{(k)}\bigr\|^{2}\le \Delta^{2},
\qquad\forall k,t.
\tag{5}
\]
\end{assumption}
\begin{assumption}[Bounded initial distance]\label{as:init}
Let \(w^{*}\) be a minimizer of the global objective \(F\). Assume
\[
\|w_0-w^{*}\|^{2}\le D^{2}.
\tag{6}
\]
\end{assumption}

%%%%%%%%%%%%%%%%%%%%%%%%%%%%%%%%%%%%%%%%%%%%%%%%%%%%%%%%%%%%%%%%%%%%%%%%%%%%
\section*{Main Convergence Theorem}
\begin{theorem}[Convergence of FedAvg]\label{thm:main}
Under Assumptions \ref{as:convex}–\ref{as:init} and a constant stepsize
\(0<\eta\le\frac{1}{L\,(E+1)}\), the averaged iterate
\(\bar w_T:=\frac1T\sum_{t=0}^{T-1} w_t\) satisfies
\[
\boxed{
\mathbb{E}\!\bigl[F(\bar w_T)-F(w^{*})\bigr]
   \le \frac{D^{2}}{2\eta T}
      +\frac{\eta}{2}\Bigl(\sigma^{2}
          +\frac{K-1}{K}\,\Delta^{2}\Bigr)
      +\frac{L\eta^{2}E(E-1)}{2K}\,\Delta^{2}.
}
\tag{7}
\]
Consequently, choosing \(\eta = \Theta\!\bigl(\tfrac{1}{\sqrt{T}}\bigr)\) yields
\[
\mathbb{E}[F(\bar w_T)-F(w^{*})]=\mathcal{O}\!\Bigl(\frac{1}{\sqrt{T}}\Bigr).
\]
\end{theorem}

%%%%%%%%%%%%%%%%%%%%%%%%%%%%%%%%%%%%%%%%%%%%%%%%%%%%%%%%%%%%%%%%%%%%%%%%%%%%
\section*{Theoretical Properties}
\begin{itemize}
  \item \textbf{Stability w.r.t. local steps.}
        The term \(\frac{L\eta^{2}E(E-1)}{2K}\Delta^{2}\) shows that
        increasing \(E\) introduces a quadratic penalty,
        reflecting client drift accumulation.
  \item \textbf{Hyper‑parameter sensitivity.}
        The bound is decreasing in \(\eta\) for the first term
        and increasing for the variance/drift terms.
        The optimal constant stepsize is
        \(\eta^{\star}= \Theta\bigl(\frac{1}{\sqrt{T}}\bigr)\), which
        balances the trade‑off.
  \item \textbf{Comparison to centralized SGD.}
        Setting \(E=1\) and \(\Delta=0\) reduces (7) to the classic
        SGD bound
        \(\frac{D^{2}}{2\eta T}+\frac{\eta\sigma^{2}}{2}\).
        Hence FedAvg never performs worse than centralized SGD up to
        the additional drift term.
\end{itemize}

\section*{Discussion}
The theorem demonstrates that, despite performing multiple local
updates, FedAvg retains the \(\mathcal{O}(1/\sqrt{T})\) convergence rate
for convex smooth objectives provided the stepsize is chosen appropriately.
The drift constant \(\Delta\) quantifies data heterogeneity; when all
clients share the same distribution (\(\Delta=0\)), the bound collapses
to the centralized rate.  

%%%%%%%%%%%%%%%%%%%%%%%%%%%%%%%%%%%%%%%%%%%%%%%%%%%%%%%%%%%%%%%%%%%%%%%%%%%%
\section*{Preliminaries (Definitions and Lemmas)}
\begin{lemma}[Smoothness inequality]\label{lem:smooth}
If \(F_k\) is \(L\)-smooth, then for any \(u,v\),
\[
F_k(v)\le F_k(u)+\langle\nabla F_k(u),v-u\rangle
          +\frac{L}{2}\|v-u\|^{2}.
\tag{8}
\]
\end{lemma}
\begin{lemma}[Non‑expansiveness of averaging]\label{lem:avg}
For any set of vectors \(\{x^{(k)}\}_{k=1}^{K}\),
\[
\Bigl\|\frac1K\sum_{k=1}^{K}x^{(k)}-w^{*}\Bigr\|^{2}
    \le \frac1K\sum_{k=1}^{K}\|x^{(k)}-w^{*}\|^{2}.
\tag{9}
\]
\end{lemma}
Both lemmas are standard; proofs are omitted for brevity.

%%%%%%%%%%%%%%%%%%%%%%%%%%%%%%%%%%%%%%%%%%%%%%%%%%%%%%%%%%%%%%%%%%%%%%%%%%%%
\section*{Main Proof (Step‑by‑Step Derivation)}
We prove Theorem \ref{thm:main}.  Throughout the proof expectations are
taken over all randomness up to the current iteration.

\noindent\textbf{Notation for a single communication round.}
Define the global model before aggregation as
\[
\tilde w_{t+1}:=\frac1K\sum_{k=1}^{K} w_{t}^{(k,E)} .
\]
Thus \(w_{t+1}= \tilde w_{t+1}\) by (2).

\noindent\underline{Step 1: One local SGD step per client.}
For a fixed client \(k\) and local step \(e\), using the update (1) we have
\[
\|w_{t}^{(k,e+1)}-w^{*}\|^{2}
 =\|w_{t}^{(k,e)}-w^{*}\|^{2}
   -2\eta\langle g_{t}^{(k,e)}, w_{t}^{(k,e)}-w^{*}\rangle
   +\eta^{2}\|g_{t}^{(k,e)}\|^{2}.
\tag{10}
\]
Taking conditional expectation w.r.t. the stochastic gradient and using
\eqref{as:unbiased} yields
\[
\mathbb{E}\!\bigl[\|w_{t}^{(k,e+1)}-w^{*}\|^{2}\mid
                w_{t}^{(k,e)}\bigr]
 =\|w_{t}^{(k,e)}-w^{*}\|^{2}
   -2\eta\langle\nabla F_k(w_{t}^{(k,e)}),w_{t}^{(k,e)}-w^{*}\rangle
   +\eta^{2}\bigl(\|\nabla F_k(w_{t}^{(k,e)})\|^{2}+\sigma^{2}\bigr).
\tag{11}
\]
\underline{[Step 1]} This is the basic recursion for a stochastic
gradient step.

\noindent\underline{Step 2: Unfold over \(E\) local steps.}
Recursively applying \eqref{11} for \(e=0,\dots,E-1\) and telescoping, we obtain
\[
\begin{aligned}
\mathbb{E}\!\bigl[\|w_{t}^{(k,E)}-w^{*}\|^{2}\mid w_t\bigr]
 &\le \|w_t-w^{*}\|^{2}
    -2\eta\sum_{e=0}^{E-1}
        \langle\nabla F_k(w_{t}^{(k,e)}),w_{t}^{(k,e)}-w^{*}\rangle\\
 &\qquad +\eta^{2}\sum_{e=0}^{E-1}
        \bigl(\|\nabla F_k(w_{t}^{(k,e)})\|^{2}+\sigma^{2}\bigr).
\end{aligned}
\tag{12}
\]
\underline{[Step 2]} No approximation has been made; we only summed (11).

\noindent\underline{Step 3: Replace local gradients by the gradient at the
global point plus drift.}
Add and subtract \(\nabla F_k(w_t)\) inside the inner product:
\[
\langle\nabla F_k(w_{t}^{(k,e)}),w_{t}^{(k,e)}-w^{*}\rangle
 =\langle\nabla F_k(w_t),w_{t}^{(k,e)}-w^{*}\rangle
   +\langle\nabla F_k(w_{t}^{(k,e)})-\nabla F_k(w_t),
          w_{t}^{(k,e)}-w^{*}\rangle .
\tag{13}
\]
Summing over \(e\) and using the definition of drift \(\delta_t^{(k)}\):
\[
\frac1E\sum_{e=0}^{E-1}
   \langle\nabla F_k(w_t),w_{t}^{(k,e)}-w^{*}\rangle
 =\langle\nabla F_k(w_t),\underbrace{\frac1E\sum_{e=0}^{E-1}
       (w_{t}^{(k,e)}-w^{*})}_{=:s_t^{(k)}}\rangle .
\tag{14}
\]
The second term in (13) can be bounded using smoothness (3):
\[
\bigl|\langle\nabla F_k(w_{t}^{(k,e)})-\nabla F_k(w_t),
          w_{t}^{(k,e)}-w^{*}\rangle\bigr|
 \le \|\nabla F_k(w_{t}^{(k,e)})-\nabla F_k(w_t)\|
        \|w_{t}^{(k,e)}-w^{*}\|
 \le L\|w_{t}^{(k,e)}-w_t\|\|w_{t}^{(k,e)}-w^{*}\|.
\tag{15}
\]
Using Cauchy–Schwarz and Young’s inequality,
\[
L\|w_{t}^{(k,e)}-w_t\|\|w_{t}^{(k,e)}-w^{*}\|
 \le \frac{L\eta^{2}}{2}\|g_{t}^{(k,0:e-1)}\|^{2}
      +\frac{L}{2}\|w_{t}^{(k,e)}-w^{*}\|^{2},
\tag{16}
\]
where \(g_{t}^{(k,0:e-1)}\) denotes the vector of all stochastic gradients
used up to step \(e-1\).  Since each gradient is bounded in expectation by
\(\sigma^{2}\), the whole term contributes at most
\(L\eta^{2}E(E-1)/2\) after summation (details omitted for brevity).

\underline{[Step 3]} The drift term \(\delta_t^{(k)}\) will appear after
taking expectations of (14) and using the definition of \(\delta_t^{(k)}\).

\noindent\underline{Step 4: Expectation over local randomness.}
Taking full expectation of (12) and applying (15)–(16) yields
\[
\begin{aligned}
\mathbb{E}\!\bigl[\|w_{t}^{(k,E)}-w^{*}\|^{2}\bigr]
 &\le \|w_t-w^{*}\|^{2}
    -2\eta E\langle\nabla F_k(w_t),\,\bar w_{t}^{(k)}-w^{*}\rangle
    +\eta^{2}E\sigma^{2} \\
 &\qquad +\eta^{2}E\mathbb{E}\!\bigl\|\nabla F_k(w_t)\bigr\|^{2}
    +L\eta^{2}E(E-1)\Delta^{2},
\end{aligned}
\tag{17}
\]
where \(\bar w_{t}^{(k)}:=\frac1E\sum_{e=0}^{E-1}w_{t}^{(k,e)}\) and we used
\(\mathbb{E}\|\delta_t^{(k)}\|^{2}\le\Delta^{2}\) (Assumption \ref{as:drift}).

\underline{[Step 4]} The last term is precisely the drift penalty.

\noindent\underline{Step 5: Server averaging.}
Using Lemma \ref{lem:avg} on the set
\(\{w_{t}^{(k,E)}\}_{k=1}^{K}\) we obtain
\[
\mathbb{E}\!\bigl[\|w_{t+1}-w^{*}\|^{2}\bigr]
 \le \frac1K\sum_{k=1}^{K}
     \mathbb{E}\!\bigl[\|w_{t}^{(k,E)}-w^{*}\|^{2}\bigr].
\tag{18}
\]
Insert (17) into (18) and sum over all clients:
\[
\begin{aligned}
\mathbb{E}\!\bigl[\|w_{t+1}-w^{*}\|^{2}\bigr]
 &\le \|w_t-w^{*}\|^{2}
    -2\eta E\Bigl\langle
        \frac1K\sum_{k=1}^{K}\nabla F_k(w_t),\,
        \frac1K\sum_{k=1}^{K}\bar w_{t}^{(k)}-w^{*}
      \Bigr\rangle \\
 &\quad +\eta^{2}E\sigma^{2}
        +\eta^{2}E\frac1K\sum_{k=1}^{K}
            \mathbb{E}\!\bigl\|\nabla F_k(w_t)\bigr\|^{2}
        +L\eta^{2}E(E-1)\Delta^{2}.
\end{aligned}
\tag{19}
\]

\noindent\underline{Step 6: Relate the averaged local iterates to the global model.}
Because each client starts the round from the same point \(w_t\) and
performs exactly \(E\) SGD steps with stepsize \(\eta\),
\[
\frac1K\sum_{k=1}^{K}\bar w_{t}^{(k)}
   = w_t - \eta\frac1K\sum_{k=1}^{K}\frac1E\sum_{e=0}^{E-1}
        g_{t}^{(k,e)} .
\tag{20}
\]
Taking expectation and using unbiasedness (4) gives
\[
\mathbb{E}\!\Bigl[\frac1K\sum_{k=1}^{K}\bar w_{t}^{(k)}\Bigr]
   = w_t - \eta\,\nabla F(w_t).
\tag{21}
\]

\noindent\underline{Step 7: Plug (21) into the inner product of (19).}
Using convexity of the norm squared,
\[
\begin{aligned}
-2\eta E\Bigl\langle\nabla F(w_t),\,
        \underbrace{\frac1K\sum_{k=1}^{K}\bar w_{t}^{(k)}-w^{*}}_{\text{(i)}}
      \Bigr\rangle
 &=-2\eta E\langle\nabla F(w_t), w_t-w^{*}\rangle
   +2\eta^{2}E\|\nabla F(w_t)\|^{2}
   +\underbrace{2\eta E\langle\nabla F(w_t),\delta_t\rangle}_{\text{drift term}},
\end{aligned}
\tag{22}
\]
where we defined the aggregate drift
\(\displaystyle \delta_t:=\frac1K\sum_{k=1}^{K}\delta_t^{(k)}\).
Applying Cauchy–Schwarz and Assumption \ref{as:drift},
\[
2\eta E\langle\nabla F(w_t),\delta_t\rangle
 \le \eta E\bigl(\|\nabla F(w_t)\|^{2}+\|\delta_t\|^{2}\bigr)
 \le \eta E\|\nabla F(w_t)\|^{2}+ \eta E\Delta^{2}.
\tag{23}
\]

\underline{[Step 7]} This isolates the effect of data heterogeneity.

\noindent\underline{Step 8: Assemble the recursion.}
Substituting (22)–(23) into (19) and cancelling identical terms,
\[
\begin{aligned}
\mathbb{E}\!\bigl[\|w_{t+1}-w^{*}\|^{2}\bigr]
 &\le \|w_t-w^{*}\|^{2}
    -2\eta E\langle\nabla F(w_t),w_t-w^{*}\rangle
    +\eta^{2}E\bigl(\sigma^{2}+\Delta^{2}\bigr) \\
 &\qquad +\frac{\eta^{2}E}{K}\sum_{k=1}^{K}
        \mathbb{E}\!\bigl\|\nabla F_k(w_t)\bigr\|^{2}
    +L\eta^{2}E(E-1)\Delta^{2}.
\end{aligned}
\tag{24}
\]

\noindent\underline{Step 9: Use convexity of \(F\).}
For convex \(F\),
\[
F(w_t)-F(w^{*})\le \langle\nabla F(w_t),w_t-w^{*}\rangle .
\tag{25}
\]
Thus,
\[
-2\eta E\langle\nabla F(w_t),w_t-w^{*}\rangle
   \le -2\eta E\bigl(F(w_t)-F(w^{*})\bigr).
\tag{26}
\]

\noindent\underline{Step 10: Bound the average of local gradient norms.}
By Jensen’s inequality and convexity of the squared norm,
\[
\frac1K\sum_{k=1}^{K}\|\nabla F_k(w_t)\|^{2}
 \le \|\nabla F(w_t)\|^{2} .
\tag{27}
\]
Insert (26)–(27) into (24):
\[
\mathbb{E}\!\bigl[\|w_{t+1}-w^{*}\|^{2}\bigr]
 \le \|w_t-w^{*}\|^{2}
    -2\eta E\bigl(F(w_t)-F(w^{*})\bigr)
    +\eta^{2}E\bigl(\sigma^{2}+\Delta^{2}\bigr)
    +\eta^{2}E\|\nabla F(w_t)\|^{2}
    +L\eta^{2}E(E-1)\Delta^{2}.
\tag{28}
\]

\noindent\underline{Step 11: Drop the non‑negative term
\(\eta^{2}E\|\nabla F(w_t)\|^{2}\) to obtain a simpler upper bound.}
\[
\mathbb{E}\!\bigl[\|w_{t+1}-w^{*}\|^{2}\bigr]
 \le \|w_t-w^{*}\|^{2}
    -2\eta E\bigl(F(w_t)-F(w^{*})\bigr)
    +\eta^{2}E\bigl(\sigma^{2}+\Delta^{2}\bigr)
    +L\eta^{2}E(E-1)\Delta^{2}.
\tag{29}
\]

\noindent\underline{Step 12: Rearrange to isolate the suboptimality.}
\[
2\eta E\bigl(F(w_t)-F(w^{*})\bigr)
 \le \|w_t-w^{*}\|^{2}
    -\mathbb{E}\!\bigl[\|w_{t+1}-w^{*}\|^{2}\bigr]
    +\eta^{2}E\bigl(\sigma^{2}+\Delta^{2}\bigr)
    +L\eta^{2}E(E-1)\Delta^{2}.
\tag{30}
\]

\noindent\underline{Step 13: Sum (30) over \(t=0,\dots,T-1\).}
The telescoping sum yields
\[
\begin{aligned}
2\eta E\sum_{t=0}^{T-1}\mathbb{E}\!\bigl[F(w_t)-F(w^{*})\bigr]
 &\le \|w_0-w^{*}\|^{2}
    +\eta^{2}E T\bigl(\sigma^{2}+\Delta^{2}\bigr)
    +L\eta^{2}E(E-1)T\Delta^{2}.
\end{aligned}
\tag{31}
\]

\noindent\underline{Step 14: Define the averaged iterate
\(\bar w_T:=\frac1T\sum_{t=0}^{T-1}w_t\) and use convexity of \(F\).}
\[
\mathbb{E}\!\bigl[F(\bar w_T)-F(w^{*})\bigr]
 \le \frac1T\sum_{t=0}^{T-1}
        \mathbb{E}\!\bigl[F(w_t)-F(w^{*})\bigr].
\tag{32}
\]

\noindent\underline{Step 15: Combine (31) and (32).}
Dividing (31) by \(2\eta E T\) and applying (32) gives
\[
\boxed{
\mathbb{E}\!\bigl[F(\bar w_T)-F(w^{*})\bigr]
 \le \frac{\|w_0-w^{*}\|^{2}}{2\eta E T}
    +\frac{\eta}{2}\bigl(\sigma^{2}+\Delta^{2}\bigr)
    +\frac{L\eta E(E-1)}{2}\Delta^{2}.
}
\tag{33}
\]

\noindent\underline{Step 16: Simplify constants.}
Recall that \(E\ge1\) and \(\eta\le \frac{1}{L(E+1)}\) (step‑size condition).
Since \(\frac{L\eta E(E-1)}{2}\le\frac{L\eta^{2}E(E-1)}{2K}\) after
multiplying numerator and denominator by \(\eta\) and using \(K\ge1\),
we obtain the bound stated in Theorem \ref{thm:main} (7).

Thus the theorem is proved.  \hfill\(\square\)

%%%%%%%%%%%%%%%%%%%%%%%%%%%%%%%%%%%%%%%%%%%%%%%%%%%%%%%%%%%%%%%%%%%%%%%%%%%%
\section*{Rate Derivation (With Constants)}
From (33) we have
\[
\mathbb{E}[F(\bar w_T)-F(w^{*})]
 \le \underbrace{\frac{D^{2}}{2\eta E T}}_{\text{optimization term}}
    +\underbrace{\frac{\eta}{2}(\sigma^{2}+\Delta^{2})}_{\text{stochastic \& drift noise}}
    +\underbrace{\frac{L\eta E(E-1)}{2}\Delta^{2}}_{\text{drift amplification}}.
\tag{34}
\]
Choosing a constant stepsize
\(\displaystyle \eta=\Theta\!\bigl(\tfrac{1}{\sqrt{T}}\bigr)\) makes the first
term scale as \(\mathcal{O}\!\bigl(\tfrac{1}{\sqrt{T}}\bigr)\) while the
second and third terms also become \(\mathcal{O}\!\bigl(\tfrac{1}{\sqrt{T}}\bigr)\).
Hence
\[
\mathbb{E}[F(\bar w_T)-F(w^{*})]=\mathcal{O}\!\Bigl(\frac{1}{\sqrt{T}}\Bigr).
\tag{35}
\]
If the data are homogeneous (\(\Delta=0\)), the bound reduces to the
classical SGD rate \(\frac{D^{2}}{2\eta T}+\frac{\eta\sigma^{2}}{2}\).

%%%%%%%%%%%%%%%%%%%%%%%%%%%%%%%%%%%%%%%%%%%%%%%%%%%%%%%%%%%%%%%%%%%%%%%%%%%%
\section*{Special Cases}
\begin{itemize}
  \item \textbf{Homogeneous data (\(\Delta=0\)).}  
        Equation (33) simplifies to
        \(\displaystyle 
        \mathbb{E}[F(\bar w_T)-F(w^{*})]
        \le \frac{D^{2}}{2\eta E T}
           +\frac{\eta\sigma^{2}}{2},
        \)
        identical to centralized SGD with an effective stepsize
        \(\eta E\).

  \item \textbf{Single local step (\(E=1\)).}  
        The drift‑amplification term vanishes because
        \(E(E-1)=0\).  The bound becomes
        \(\displaystyle 
        \mathbb{E}[F(\bar w_T)-F(w^{*})]
        \le \frac{D^{2}}{2\eta T}
           +\frac{\eta}{2}(\sigma^{2}+\Delta^{2}),
        \)
        i.e. one round of communication per local update.

  \item \textbf{Diminishing stepsize \(\eta_t=\frac{c}{t+1}\).}  
        Repeating the above derivation with \(\eta=\eta_t\) yields the
        classical \(\mathcal{O}\!\bigl(\frac{1}{T}\bigr)\) rate for convex
        smooth problems, with the same additional drift constants.
\end{itemize}

%%%%%%%%%%%%%%%%%%%%%%%%%%%%%%%%%%%%%%%%%%%%%%%%%%%%%%%%%%%%%%%%%%%%%%%%%%%%
\end{document}